% ===== Chapter 5 opening: perception is a reconstruction (~4 min) =====
\partopener{Գլուխ 5}{Ընկալում}{Կանխատեսե՞լ էիր այս գլուխը. ինչպես է ուղեղը ընկալում աշխարհը (ու որտեղ է սխալվում)}{confused}{peach}

\begin{frame}{Ներս նայել ու դուրս նայել}
  \begin{columns}[T]
    \begin{column}{0.47\textwidth}
      \begin{card}[colback=lavender!40!cardwhite]
        \emoji[1.4em]{thought}\ {\blackfont Ներս}\par\smallskip
        Մենք ինքնագիտակից ենք. կարող ենք զգալ ու ուսումնասիրել սեփական միտքը (introspection)։ Շատերը սա են ամենաշատը գնահատում։
      \end{card}
    \end{column}
    \begin{column}{0.47\textwidth}
      \onslide<2->{%
      \begin{card}[colback=peach!40!cardwhite]
        \emoji[1.4em]{eyes}\ {\blackfont Դուրս}\par\smallskip
        Բայց ուղեղի մեխանիզմների մեծ մասը զբաղված է նրանով, թե ինչպես ընկալել աշխարհը \textbf{գանգից դուրս}։
      \end{card}}
    \end{column}
  \end{columns}
  \vspace{4mm}
  \onslide<3->{%
  \centering
  \begin{tikzpicture}
    \node[wbox=peach!40!cardwhite, text width=3cm] (a) at (0,0) {\emoji{nose}\emoji{ear}\emoji{eye}\ ընկալում ենք};
    \node[wbox=butter!50!cardwhite, text width=3cm] (b) at (4.4,0) {\emoji{target}\ ընտրում ենք կարևորը};
    \node[wbox=mint!45!cardwhite, text width=3cm] (c) at (8.8,0) {\emoji{runner}\ գործում ենք};
    \draw[warr] (a) -- (b);
    \draw[warr] (b) -- (c);
  \end{tikzpicture}\par}
\end{frame}

\begin{frame}{Առասպել, թե՞ փաստ}
  \mythorfact{Աչքերն ու ականջները տեսախցիկ ու խոսափող են, ուղեղը միայն դիտում է}{ԱՌԱՍՊԵԼ}{pinkdark}{%
    \vspace{1mm}
    \begin{card}[colback=brainpink!25!cardwhite]
      Զգայարաններից եկող հում տվյալները ոչ թե հարուստ հոսք են, այլ \textbf{պղտոր առվակ}։ Ուղեղն է դրանից «հղկում» աշխարհի մեր հարուստ պատկերը։ Կենսաբանությունը տեխնոլոգիա չէ։
    \end{card}}
\end{frame}

% Burnett's police sketch artist
\begin{frame}{Ուղեղը՝ ոստիկանական նկարիչ}
  \centering
  % Burnett's analogy: the brain as a police sketch artist who gets hundreds of descriptions at once
% and must render a full-colour 3D town, updated every minute. Reveals in 3 steps (use inside a frame).
\begin{tikzpicture}[
    witness/.style={draw=plum, line width=0.7pt, rounded corners=2.5mm, inner sep=2.5pt,
      font=\footnotesize, text=plum, anchor=west},
  ]
  \path[use as bounding box] (-6.9,-1.85) rectangle (6.9,1.75);
  % --- witnesses (the senses) -------------------------------------------
  \foreach \y/\e/\t/\c/\x in {1.36/eye/{«Կարմիր մի բան էր»}/peach/-6.8, 0.68/ear/{«Ձայնը ձախից էր»}/mint/-6.5,
      0/nose/{«Հացի հոտ էր գալիս»}/butter/-6.8, -0.68/hand/{«Սառն էր»}/lavender/-6.4,
      -1.36/tongue/{«Աղի՞ էր...»}/brainpink/-6.7}{
    \node[witness, fill=\c!45!cardwhite] (w\e) at (\x,\y) {\emoji[1em]{\e}\ \t};
    \draw[-{Stealth[length=2mm]}, plum!70, line width=0.8pt] (w\e.east) -- (-2.2,0.1*\y);}
  \node[font=\mediumfont\scriptsize, text=plum!75!cream, visible on=<3->, anchor=west] at (-6.8,1.78)
    {...ու ևս հարյուրավորը, միաժամանակ};
  % --- the artist -----------------------------------------------------------
  \node at (-1.1,0.15) {\mascot[1.9cm]{thinking}};
  \node[rotate=20] at (-0.25,-0.6) {\bigemoji[0.65cm]{pencil}};
  \node[font=\blackfont\small, text=plum] at (-1.1,-1.4) {նկարիչը};
  % --- the drawing: a whole town in 3D ----------------------------------
  \begin{scope}[visible on=<2->]
    \draw[warr] (0.25,0.1) -- (1.05,0.1);
    \filldraw[fill=cardwhite, draw=plum, line width=0.9pt, rounded corners=2mm] (1.2,-1.6) rectangle (6.8,1.6);
    \fill[mint!40!cardwhite, rounded corners=2mm] (1.25,-1.55) rectangle (6.75,-0.6);
    \fill[amOrange!70] (6.15,1.05) circle (0.26);
    \foreach \x/\h/\c in {1.7/0.95/peach, 2.75/1.4/lavender, 3.9/0.85/brainpink, 4.9/1.2/butter}{
      \filldraw[fill=\c, draw=plum, line width=0.7pt] (\x,-0.9) rectangle ++(0.8,\h);
      \filldraw[fill=plum!70, draw=plum, line width=0.7pt] (\x-0.08,-0.9+\h) -- (\x+0.4,-0.9+\h+0.38) -- (\x+0.88,-0.9+\h) -- cycle;
      \fill[cardwhite] (\x+0.15,-0.9+\h-0.4) rectangle ++(0.2,0.2);
      \fill[cardwhite] (\x+0.45,-0.9+\h-0.4) rectangle ++(0.2,0.2);}
    \node at (6.15,-1.05) {\bigemoji[0.55cm]{tree}};
    \node at (2.35,-1.2) {\bigemoji[0.45cm]{walking}};
    \node at (3.75,-1.2) {\bigemoji[0.45cm]{car}};
    \node[font=\blackfont\scriptsize, text=plum, anchor=north west] at (1.35,1.53) {գունավոր, 3D, բոլոր բնակիչներով};
  \end{scope}
  \node[visible on=<3->, draw=plum, line width=0.8pt, fill=butter, rounded corners=2mm, rotate=-6,
        font=\blackfont\scriptsize, text=plum, inner sep=3pt] at (5.5,-1.65) {\emoji[1em]{alarm}\ թարմացնել ամեն րոպե};
\end{tikzpicture}
\par
  \onslide<4->{%
  \begin{bookquote}
    Ուղեղը մոտավորապես այդպես է աշխատում. միայն թե, հավանաբար, այդ նկարչի պես ջղային չէ։
  \end{bookquote}}
\end{frame}

% Overview map: where each sense arrives
\begin{frame}{Ո՞ւր են գնում զգայարանների ազդանշանները}
  \begin{columns}[c]
    \begin{column}{0.6\textwidth}
      \centering
      % Where the senses arrive in the cortex: smell (olfactory bulb), hearing (auditory cortex, temporal lobe),
% touch (somatosensory cortex, parietal lobe), vision (visual cortex, occipital lobe). Reveals 1-4.
\begin{tikzpicture}[
    tag/.style={draw=plum, line width=0.7pt, fill=cardwhite, rounded corners=2mm, inner sep=2.2pt,
      font=\scriptsize, text=plum, align=center},
    pin/.style={plum, line width=0.8pt},
  ]
  \path[use as bounding box] (-4.3,-2.45) rectangle (4.3,2.45);
  \begin{scope}[scale=0.74, yshift=0.15cm]
    \brainlobes{lavender!55!cream}{mint!55!cream}{peach!60!cream}{butter!70!cream}
    % smell: olfactory bulb under the frontal lobe
    \filldraw[fill=pinkdark, draw=plum, line width=0.6pt, rotate around={-8:(-2.45,-0.78)}] (-2.45,-0.78) ellipse (0.5 and 0.13);
    % hearing: auditory cortex on the upper temporal lobe
    \fill[amOrange, opacity=0.75, visible on=<2->] (0.35,-0.2) ellipse (0.75 and 0.24);
    % touch: somatosensory strip behind the central sulcus
    \draw[amBlue, opacity=0.55, line width=5pt, line cap=round, visible on=<3->]
      plot[smooth] coordinates {(0.55,2.55) (0.33,1.85) (0.2,1.15) (-0.12,0.4)};
    % vision: visual cortex at the back
    \fill[amRed, opacity=0.6, visible on=<4->] (3.0,-0.15) ellipse (0.32 and 0.6);
    % labels
    \node[tag] (ts) at (-4.0,-2.6) {\emoji[1em]{nose}\ հոտառական\\ սոխուկ};
    \draw[pin] (ts.north) -- (-2.6,-0.85);
    \node[tag, visible on=<2->] (th) at (0.2,-2.75) {\emoji[1em]{ear}\ լսողական կեղև\\ (քունքային բիլթ)};
    \draw[pin, visible on=<2->] (th.north) -- (0.35,-0.3);
    \node[tag, visible on=<3->] (tt) at (-2.6,2.95) {\emoji[1em]{hand}\ զգացողական կեղև\\ (գագաթային բիլթ)};
    \draw[pin, visible on=<3->] (tt.east) -- (0.45,2.3);
    \node[tag, visible on=<4->] (tv) at (4.2,-2.6) {\emoji[1em]{eye}\ տեսողական կեղև\\ (ծոծրակային բիլթ)};
    \draw[pin, visible on=<4->] (tv.north) -- (3.05,-0.65);
  \end{scope}
\end{tikzpicture}

    \end{column}
    \begin{column}{0.36\textwidth}
      \begin{itemize}
        \item Զգայարանը միջավայրից ինչ-որ բան է որսում ու դարձնում \textbf{էլեկտրաքիմիական ազդանշան}
        \item<5-> Ամեն ազդանշան իր «հասցեն» ունի ուղեղում
      \end{itemize}
      \onslide<5->{%
      \begin{card}[colback=butter!45!cream]
        \small Քիչ տվյալից մանրամասն պատկեր սարքելիս սխալներն անխուսափելի են։
      \end{card}}
    \end{column}
  \end{columns}
\end{frame}
